% Part: first-order-logic
% Chapter: natural-deduction
% Section: proof-theoretic-notions

\documentclass[../../../include/open-logic-section]{subfiles}

\begin{document}

\iftag{FOL}
      {\olfileid[fr]{fol}{ntd}{ptn}}
      {\olfileid[fr]{pl}{ntd}{ptn}}

\Needspace{9\baselineskip}
\olsection{Notions de théorie de la démonstration}

\begin{editorial}
Cette section rassemble les définitions de la relation de
!!{derivability} et de la cohérence pour la déduction naturelle.
\end{editorial}

\begin{explain}
Nous avons défini plusieurs notions sémantiques importantes :
validité, conséquence logique et satisfaisabilité. Définissons
maintenant les \emph{notions de théorie de la démonstration}
correspondantes. Elles ne font pas appel à la satisfaction
d'!!{sentence}s dans des \iftag{FOL}{!!{structure}s}{!!{valuation}s},
mais à la !!{derivability} ou à la !!{nonderivability} de certains
!!{sentence}s à partir d'autres. La coïncidence de ces notions
constitue une découverte importante : c'est le contenu des
\emph{théorèmes de correction et de complétude}.
\end{explain}

\begin{defn}[Théorèmes]
!!^a{sentence}~$!A$ est un \emph{théorème} s'il existe une
!!{derivation} de~$!A$ en déduction naturelle dans laquelle toutes
les hypothèses sont !!{discharged}s. On écrit $\Proves !A$ si
$!A$ est un théorème, et $\Proves/ !A$ dans le cas contraire.
\end{defn}

\begin{defn}[!!^{derivability}]
!!^a{sentence}~$!A$ est \emph{!!{derivable} à partir} d'un ensemble
d'!!{sentence}s~$\Gamma$, ce que l'on note $\Gamma \Proves !A$,
s'il existe une !!{derivation} de conclusion~$!A$ dont chaque
hypothèse est soit !!{discharged}, soit un élément de~$\Gamma$.
Si $!A$ n'est pas !!{derivable} à partir de $\Gamma$, on écrit
$\Gamma \Proves/ !A$.
\end{defn}

\begin{defn}[Cohérence]
Un ensemble d'!!{sentence}s~$\Gamma$ est \emph{incohérent} si et
seulement si $\Gamma \Proves \lfalse$. S'il n'est pas incohérent,
c'est-à-dire si $\Gamma \Proves/ \lfalse$, on dit qu'il est
\emph{cohérent}.
\end{defn}

\begin{prop}[Réflexivité]
\ollabel{prop:reflexivity}
Si $!A \in \Gamma$, alors $\Gamma \Proves !A$.
\end{prop}

\begin{proof}
L'hypothèse $!A$ seule constitue une !!{derivation} de~$!A$ dont
toutes les hypothèses !!{undischarged}s, en l'occurrence $!A$,
appartiennent à~$\Gamma$.
\end{proof}

\begin{prop}[Monotonie]
\ollabel{prop:monotonicity}
Si $\Gamma \subseteq \Delta$ et $\Gamma \Proves !A$, alors
$\Delta \Proves !A$.
\end{prop}

\begin{proof}
Toute !!{derivation} de $!A$ à partir de $\Gamma$ est aussi une
!!{derivation} de $!A$ à partir de~$\Delta$.
\end{proof}

\begin{prop}[Transitivité]
\ollabel{prop:transitivity}
Si $\Gamma \Proves !A$ et $\{!A\} \cup \Delta \Proves !B$,
alors $\Gamma \cup \Delta \Proves !B$.
\end{prop}

\begin{proof}
Si $\Gamma \Proves !A$, il existe une !!{derivation}~$\delta_0$
de~$!A$ dont toutes les hypothèses !!{undischarged}s appartiennent
à~$\Gamma$. Si $\{!A\} \cup \Delta \Proves !B$, il existe une
!!{derivation}~$\delta_1$ de~$!B$ dont toutes les hypothèses
!!{undischarged}s appartiennent à~$\{!A\} \cup \Delta$.
Considérons alors l'arbre suivant :
\begin{prooftree}
  \AxiomC{$\Delta, \Discharge{!A}{1}$}
  \RightLabel{$\delta_1$}
  \DeduceC{$!B$}
  \DischargeRule{\Intro{\lif}}{1}
  \UnaryInfC{$!A \lif !B$}
  \AxiomC{$\Gamma$}
  \RightLabel{$\delta_0$}
  \DeduceC{$!A$}
  \RightLabel{\Elim{\lif}}
  \BinaryInfC{$!B$}
\end{prooftree}
Toutes les hypothèses !!{undischarged}s appartiennent maintenant
à $\Gamma \cup \Delta$; nous avons donc
$\Gamma \cup \Delta \Proves !B$.
\end{proof}

Lorsque $\Gamma = \{!A_1, !A_2, \ldots, !A_k\}$ est fini, on peut
abréger $\Gamma \Proves !B$ en
$!A_1,!A_2,\ldots,!A_k \Proves !B$. En particulier,
$!A \Proves !B$ signifie $\{!A\} \Proves !B$.

Remarquons que, si $\Gamma \Proves !A$ et $!A \Proves !B$, alors
$\Gamma \Proves !B$. Il en résulte aussi que, si
$!A_1, \dots, !A_n \Proves !B$ et $\Gamma \Proves !A_i$ pour
chaque~$i$, alors $\Gamma \Proves !B$.

\begin{prop}
\ollabel{prop:incons}
Les conditions suivantes sont équivalentes :
    \begin{enumerate}
      \item \( \Gamma \) est incohérent.
      \item \( \Gamma \Proves {!A} \) pour tout !!{sentence}~\( {!A} \).
      \item \( \Gamma \Proves {!A} \) et \( \Gamma \Proves \lnot {!A} \) pour au moins un !!{sentence}~\( {!A} \).
    \end{enumerate}
\end{prop}

\begin{proof}
Exercice.
\end{proof}

\tagprob{FOL}
\begin{prob}
Démontrer la \olref[fol][ntd][ptn]{prop:incons}.
\end{prob}
\tagendprob

\tagprob{notFOL}
\begin{prob}
Démontrer la \olref[pl][ntd][ptn]{prop:incons}.
\end{prob}
\tagendprob

\begin{prop}[Compacité]
  \ollabel{prop:proves-compact}
  \begin{enumerate}
  \item Si $\Gamma \Proves !A$, il existe une partie finie
    $\Gamma_0 \subseteq \Gamma$ telle que $\Gamma_0 \Proves !A$.
  \item Si toute partie finie de~$\Gamma$ est cohérente,
    alors $\Gamma$ est cohérent.
  \end{enumerate}
\end{prop}

\begin{proof}
  \begin{enumerate}
    \item Si $\Gamma \Proves !A$, il existe une
      !!{derivation}~$\delta$ de~$!A$ à partir de~$\Gamma$.
      Soit $\Gamma_0$ l'ensemble des hypothèses
      !!{undischarged}s de~$\delta$. Toute !!{derivation} étant
      finie, $\Gamma_0$ ne contient qu'un nombre fini
      d'!!{sentence}s. Ainsi, $\delta$ est une !!{derivation}
      de~$!A$ à partir d'une partie finie~$\Gamma_0 \subseteq \Gamma$.
    \item C'est la contraposée de (1) dans le cas particulier
      $!A \ident \lfalse$.
  \end{enumerate}
\end{proof}

\begin{editorial}
Dans le paragraphe introductif, la mention des structures de la
source est adaptée en valuations pour la version propositionnelle.
Les contextes d'hypothèses restent des ensembles éventuellement
infinis; seule chaque dérivation particulière est finie.
\end{editorial}

\end{document}
